\documentclass[12pt,a4paper]{article}

% ============================================================
% Drake Journal Reading LaTeX Template
% 作者: 謝慕揚醫師 MD, PhD, FESC
% 
% 使用說明:
% 1. 表格使用 longtable，不要有垂直分隔線 |
% 2. 表格內換行使用 \newline，不要使用 \\
% 3. 藥物名稱、檢驗、檢查請維持英文
% 4. Reference 格式請使用 \href 加上驗證過的 DOI
% 5. 使用 XeLaTeX 編譯
% ============================================================

% Chinese font setup
\usepackage{xeCJK}
\usepackage{fontspec}
\IfFontExistsTF{Noto Serif CJK TC}{
  \setCJKmainfont{Noto Serif CJK TC}
  \setCJKsansfont{Noto Serif CJK TC}
  \setCJKmonofont{Noto Serif CJK TC}
}{\setCJKmainfont{SimSun}
  \setCJKsansfont{SimSun}
  \setCJKmonofont{SimSun}
}

% Essential packages
\usepackage[margin=2cm]{geometry}
\usepackage{graphicx}
\usepackage{booktabs}
\usepackage{longtable}
\usepackage{array}
\usepackage{xcolor}
\usepackage{hyperref}
\usepackage{tcolorbox}
\usepackage{titlesec}
\usepackage{enumitem}
\usepackage{amsmath}
\usepackage{multirow}
\usepackage{fancyhdr}

% Color definitions
\definecolor{emergency_red}{RGB}{186,24,27}
\definecolor{emergency_blue}{RGB}{0,114,188}
\definecolor{emergency_gray}{RGB}{120,120,120}
\definecolor{highlight_yellow}{RGB}{255,250,205}

% Hyperlink setup
\hypersetup{
    colorlinks=true,
    linkcolor=emergency_blue,
    citecolor=emergency_blue,
    urlcolor=emergency_blue,
    pdfborder={0 0 0}
}

% Title formats
\titleformat{\section}{\Large\bfseries\color{emergency_red}}{\thesection}{1em}{}
\titleformat{\subsection}{\large\bfseries\color{emergency_blue}}{\thesubsection}{1em}{}
\titleformat{\subsubsection}{\normalsize\bfseries\color{emergency_gray}}{\thesubsubsection}{1em}{}

% Header/Footer setup
\pagestyle{fancy}
\fancyhf{}
\fancyhead[L]{文件標題}
\fancyhead[R]{\thepage}
\renewcommand{\headrulewidth}{0.4pt}

% ============================================================
% 文件資訊 - 請修改以下內容
% ============================================================

\title{\textbf{English Title Here\\
中文標題在此}}
\author{謝慕揚醫師 MD, PhD, FESC\\
資料來源：Journal Name (Year)}
\date{整理日期：2025年XX月XX日}

\begin{document}

\maketitle

% Key findings box
\begin{tcolorbox}[colback=yellow!10!white,colframe=orange!75!black,title=核心發現摘要]
\textbf{重大發現：}在此描述本文最重要的發現...

\textbf{臨床意義：}對臨床實務的影響...
\end{tcolorbox}

\tableofcontents
\newpage

% ============================================================
\section{研究概述}
% ============================================================

本研究為一項前瞻性多中心觀察研究，比較...（研究簡介）

% ============================================================
\section{研究資訊}
% ============================================================

\subsection{基本資料}

\textbf{標題：}研究標題\\
\textbf{作者：}First Author, Second Author, et al.\\
\textbf{機構：}研究機構名稱\\
\textbf{期刊：}Journal Name (Year)\\
\textbf{DOI：}\href{https://doi.org/10.xxxx/xxxxx}{10.xxxx/xxxxx}\\
\textbf{研究註冊：}ClinicalTrials.gov NCT00000000

\subsection{研究背景與目標}

\begin{itemize}
    \item 背景說明一
    \item 背景說明二
    \item \textbf{研究目標：}本研究旨在...
\end{itemize}

% ============================================================
\section{研究方法}
% ============================================================

\subsection{研究設計}

\begin{itemize}
    \item \textbf{研究類型：}前瞻性多中心觀察研究
    \item \textbf{研究期間：}2020-2024年
    \item \textbf{追蹤期間：}中位數 X 年
    \item \textbf{分析方法：}傾向性評分匹配
\end{itemize}

\subsection{納入與排除條件}

\begin{longtable}{p{3cm}p{11cm}}
\toprule
\textbf{類型} & \textbf{標準} \\
\midrule
\endfirsthead

\toprule
\textbf{類型} & \textbf{標準} \\
\midrule
\endhead

\bottomrule
\endfoot

\textbf{納入條件} & 
• 年齡 20-85 歲\newline
• 條件二\newline
• 條件三 \\
\midrule
\textbf{排除條件} & 
• 排除條件一\newline
• 排除條件二\newline
• 排除條件三 \\
\bottomrule
\end{longtable}

% ============================================================
\section{主要結果}
% ============================================================

\subsection{基線特徵}

\begin{longtable}{lcc}
\caption{基線患者特徵} \\
\toprule
\textbf{特徵} & \textbf{治療組 (n=XXX)} & \textbf{對照組 (n=XXX)} \\
\midrule
\endfirsthead

\multicolumn{3}{c}{\textbf{表格續接}} \\
\toprule
\textbf{特徵} & \textbf{治療組} & \textbf{對照組} \\
\midrule
\endhead

\bottomrule
\endfoot

年齡 (歲) & 65.2±10.3 & 64.8±9.8 \\
男性 & 142 (71\%) & 138 (69\%) \\
BMI (kg/m²) & 26.5±4.2 & 26.8±4.5 \\
\midrule
\multicolumn{3}{l}{\textbf{病史}} \\
糖尿病 & 68 (34\%) & 72 (36\%) \\
高血壓 & 156 (78\%) & 148 (74\%) \\
\bottomrule
\end{longtable}

\subsection{主要終點結果}

\begin{longtable}{p{5cm}ccc}
\caption{Cox比例風險迴歸分析} \\
\toprule
\textbf{終點} & \textbf{HR} & \textbf{95\% CI} & \textbf{P值} \\
\midrule
\endfirsthead

\toprule
\textbf{終點} & \textbf{HR} & \textbf{95\% CI} & \textbf{P值} \\
\midrule
\endhead

\bottomrule
\endfoot

\textbf{主要複合終點} & 0.72 & 0.52-0.99 & 0.018 \\
心血管死亡 & 0.17 & 0.07-0.40 & <0.001 \\
全因死亡 & 0.66 & 0.47-0.93 & 0.015 \\
\bottomrule
\end{longtable}

% ============================================================
\section{討論與臨床意義}
% ============================================================

\subsection{重要發現}

\begin{enumerate}
    \item \textcolor{blue}{\textbf{發現一：}}描述第一個重要發現...
    \item \textcolor{blue}{\textbf{發現二：}}描述第二個重要發現...
    \item \textcolor{blue}{\textbf{發現三：}}描述第三個重要發現...
\end{enumerate}

\subsection{臨床實務建議}

\begin{tcolorbox}[
    colback=emergency_blue!5!white,
    colframe=emergency_blue,
    title=\textbf{重點提醒},
    fonttitle=\bfseries
]
根據本研究結果，臨床建議如下：
\begin{enumerate}
    \item 建議一
    \item 建議二
    \item 建議三
\end{enumerate}
\end{tcolorbox}

% ============================================================
\section{研究限制}
% ============================================================

\begin{enumerate}
    \item \textbf{限制一：}描述限制...
    \item \textbf{限制二：}描述限制...
    \item \textbf{限制三：}描述限制...
\end{enumerate}

% ============================================================
\section{結論}
% ============================================================

\begin{enumerate}
    \item \textcolor{red}{\textbf{結論一：}}主要結論描述...
    \item \textcolor{red}{\textbf{結論二：}}次要結論描述...
\end{enumerate}

\vspace{1cm}
\textbf{關鍵詞：}關鍵詞一、關鍵詞二、關鍵詞三

% ============================================================
\section{參考文獻}
% ============================================================

\begin{enumerate}
    \item Author1 Initials, Author2 Initials, Author3 Initials, et al. Article Title. \href{https://doi.org/10.xxxx/xxxxx}{\textit{Journal Name}. Year;Volume(Issue):Pages.}
    
    \item Second Reference Author et al. Second Article Title. \href{https://doi.org/10.xxxx/xxxxx}{\textit{Journal Name}. Year;Volume:Pages.}
\end{enumerate}

% ============================================================
\section{縮寫對照表}
% ============================================================

\begin{longtable}{p{4cm}p{10cm}}
\toprule
\textbf{縮寫} & \textbf{全名} \\
\midrule
\endfirsthead

\toprule
\textbf{縮寫} & \textbf{全名} \\
\midrule
\endhead

\bottomrule
\endfoot

HR & Hazard Ratio (風險比) \\
CI & Confidence Interval (信賴區間) \\
LVEF & Left Ventricular Ejection Fraction (左室射血分數) \\
\bottomrule
\end{longtable}

\end{document}
